\documentclass[aspectratio=1610, 9pt]{beamer}
\usepackage{mobi-beamer}

\title[Proceso de Poisson]{El Proceso de Poisson}
\subtitle{Fundamentos, Distribuciones Clave y Propiedades}
\author{Alejandra Tabares Pozos}
\institute{Universidad de los Andes \\
 Departamento de Ingeniería Industrial}
\date{}

\begin{document}

% ============================== SLIDE 1 ==============================
\begin{frame}
\titlepage
\end{frame}

% ============================== SLIDE 2 ==============================
\begin{frame}{Contenido}
\tableofcontents
\end{frame}

% ======================================================================
\section{Introducci\'{o}n y Motivaci\'{o}n}
% ======================================================================

% ============================== SLIDE 3 ==============================
\begin{frame}{?`Por qu\'{e} estudiar el proceso de Poisson?}

El proceso de Poisson es el modelo fundamental para eventos que ocurren \textbf{aleatoriamente en el tiempo}.

\vspace{0.4cm}
\begin{columns}
\column{0.5\textwidth}
\textbf{Ejemplos cotidianos:}
\begin{itemize}
    \item Llamadas a un call center
    \item Llegadas de clientes a un banco
    \item Accidentes en una autopista
    \item Correos electr\'{o}nicos recibidos
    \item Fallas en un servidor
\end{itemize}

\column{0.5\textwidth}
\textbf{?`Qu\'{e} tienen en com\'{u}n?}
\begin{itemize}
    \item Ocurren uno a la vez
    \item Son ``impredecibles'' individualmente
    \item Tienen una tasa promedio estable
    \item Lo que pas\'{o} antes no afecta lo que viene
\end{itemize}
\end{columns}

\vspace{0.4cm}
\begin{block}{Objetivo}
Formalizar estas ideas intuitivas en un modelo matem\'{a}tico riguroso con distribuciones bien definidas.
\end{block}
\end{frame}

% ============================== SLIDE 4 ==============================
\begin{frame}{Ejemplo motivador: llamadas al 123}

\textbf{Situaci\'{o}n:} La l\'{i}nea de emergencias 123 recibe en promedio $\lambda = 4$ llamadas por minuto.

\vspace{0.3cm}
\begin{center}
\begin{tikzpicture}[scale=0.95]
    % Línea de tiempo
    \draw[->, thick] (0,0) -- (13,0) node[right] {$t$ (min)};
    \foreach \x in {0,1,...,12} {
        \draw (\x,0.1) -- (\x,-0.1) node[below] {\tiny \x};
    }
    % Eventos (llamadas) - marcados con flechas
    \foreach \x/\c in {0.3/1, 0.7/2, 1.1/3, 1.8/4, 2.1/5, 2.4/6, 2.9/7, 3.5/8,
        4.0/9, 4.2/10, 4.9/11, 5.6/12, 5.8/13, 6.3/14, 6.5/15, 7.1/16,
        7.8/17, 8.0/18, 8.3/19, 8.9/20, 9.2/21, 9.8/22, 10.1/23, 10.5/24,
        10.9/25, 11.2/26, 11.6/27, 12.0/28} {
        \draw[rojoclaro, thick] (\x,0) -- (\x,0.5);
        \fill[rojoclaro] (\x,0.5) circle (1.5pt);
    }
    % Anotaciones
    \draw[<->, thick, azuloscuro] (0,-0.6) -- (3,-0.6);
    \node[below, azuloscuro] at (1.5,-0.6) {\small 7 llamadas en 3 min};
    \draw[<->, thick, verdeoscuro] (5,-0.6) -- (8,-0.6);
    \node[below, verdeoscuro] at (6.5,-0.6) {\small 5 llamadas en 3 min};
\end{tikzpicture}
\end{center}

\vspace{0.2cm}
\textbf{Preguntas naturales:}
\begin{itemize}
    \item ?`Cu\'{a}l es la probabilidad de recibir exactamente 10 llamadas en 2 minutos?
    \item ?`Cu\'{a}nto tiempo pasar\'{a} hasta la pr\'{o}xima llamada?
    \item ?`Cu\'{a}nto tiempo hasta la 5\textsuperscript{a} llamada?
\end{itemize}

El proceso de Poisson responde \textbf{todas} estas preguntas.
\end{frame}

% ======================================================================
\section{Proceso de Conteo}
% ======================================================================

% ============================== SLIDE 5 ==============================
\begin{frame}{Proceso de conteo: definici\'{o}n}

\begin{block}{Definici\'{o}n: Proceso de conteo}
Un \textbf{proceso de conteo} $\{N(t), \; t \geq 0\}$ es un proceso estoc\'{a}stico donde $N(t)$ representa el n\'{u}mero total de eventos ocurridos hasta el tiempo $t$.
\end{block}

\vspace{0.3cm}
\textbf{Propiedades b\'{a}sicas:}
\begin{enumerate}
    \item $N(t) \geq 0$ para todo $t$
    \item $N(t)$ toma valores enteros
    \item Si $s < t$, entonces $N(s) \leq N(t)$
    \item $N(t) - N(s)$ = n\'{u}mero de eventos en el intervalo $(s, t]$
\end{enumerate}

\vspace{0.3cm}
\textbf{Ejemplo (123):} Si a las 10:00 van 15 llamadas y a las 10:03 van 27:
\[
N(3) - N(0) = 27 - 15 = 12 \text{ llamadas en esos 3 minutos}
\]
\end{frame}

% ============================== SLIDE 6 ==============================
\begin{frame}{Visualizaci\'{o}n: trayectoria de $N(t)$}

\begin{center}
\begin{tikzpicture}
\begin{axis}[
    width=13cm, height=6.5cm,
    xlabel={Tiempo $t$ (minutos)},
    ylabel={$N(t)$: n\'{u}mero acumulado de llamadas},
    xmin=0, xmax=5,
    ymin=0, ymax=22,
    grid=both, grid style={gray!20},
    axis lines=left,
    const plot mark right,
    thick,
    ytick={0,2,...,22},
]
\addplot[azuloscuro, line width=1.5pt, const plot mark right] coordinates {
    (0,0) (0.15,1) (0.4,2) (0.55,3) (0.9,4) (1.05,5) (1.2,6) (1.5,7) (1.75,8)
    (2.0,9) (2.1,10) (2.45,11) (2.8,12) (2.9,13) (3.15,14) (3.25,15) (3.55,16)
    (3.9,17) (4.0,18) (4.15,19) (4.45,20) (5.0,21)
};

% Incrementos
\draw[<->, thick, rojoclaro] (axis cs:1,0.5) -- (axis cs:2,0.5);
\node[below, rojoclaro, font=\small] at (axis cs:1.5,0.5) {$N(2)-N(1)=4$};

\draw[<->, thick, verdeoscuro] (axis cs:3,0.5) -- (axis cs:4,0.5);
\node[below, verdeoscuro, font=\small] at (axis cs:3.5,0.5) {$N(4)-N(3)=3$};

\end{axis}
\end{tikzpicture}
\end{center}

Funci\'{o}n escalonada creciente: salta en +1 cada vez que llega una llamada.
\end{frame}

% ======================================================================
\section{Definici\'{o}n del Proceso de Poisson}
% ======================================================================

% ============================== SLIDE 7 ==============================
\begin{frame}{Axiomas del proceso de Poisson}

\begin{block}{Definici\'{o}n: Proceso de Poisson con tasa $\lambda$}
$\{N(t), t \geq 0\}$ es un proceso de Poisson con tasa $\lambda > 0$ si:
\end{block}

\vspace{0.2cm}
\begin{enumerate}
    \item \textbf{Inicio en cero:} $N(0) = 0$.

    \vspace{0.2cm}
    \item \textbf{Incrementos independientes:} El n\'{u}mero de eventos en intervalos disjuntos son variables aleatorias independientes.

    \vspace{0.2cm}
    \item \textbf{Incrementos estacionarios:} La distribuci\'{o}n de $N(t+s) - N(s)$ depende solo de la longitud $t$, no de la posici\'{o}n $s$.

    \vspace{0.2cm}
    \item \textbf{Sin eventos simult\'{a}neos:} Para $h$ peque\~{n}o:
    \begin{align}
        P(N(h) = 1) &= \lambda h + o(h) \\
        P(N(h) \geq 2) &= o(h)
    \end{align}
    donde $o(h)/h \to 0$ cuando $h \to 0$.
\end{enumerate}
\end{frame}

% ============================== SLIDE 8 ==============================
\begin{frame}{Interpretaci\'{o}n intuitiva de los axiomas}

\textbf{Ejemplo:} llamadas al 123 con $\lambda = 4$ llamadas/min.

\vspace{0.3cm}
\begin{columns}
\column{0.5\textwidth}
\textbf{Incrementos independientes:}

Las llamadas entre 10:00--10:05 no afectan las de 10:05--10:10.

\vspace{0.2cm}
Que hayan llegado muchas llamadas en los \'{u}ltimos 5 min \textbf{no cambia} la probabilidad de llamadas futuras.

\vspace{0.3cm}
\textbf{Incrementos estacionarios:}

La probabilidad de recibir 3 llamadas en 1 minuto es la misma a las 10:00 que a las 15:00.

\column{0.5\textwidth}
\textbf{Sin eventos simult\'{a}neos:}

En un instante infinitesimal $dt$:
\begin{itemize}
    \item Prob.\ de 1 llamada $\approx 4\,dt$
    \item Prob.\ de 2+ llamadas $\approx 0$
    \item Prob.\ de 0 llamadas $\approx 1 - 4\,dt$
\end{itemize}

\vspace{0.3cm}
Si $dt = 0.001$ min (0.06 seg):
\begin{itemize}
    \item $P(\text{1 llamada}) \approx 0.004$
    \item $P(\text{0 llamadas}) \approx 0.996$
    \item $P(\text{2+ llamadas}) \approx 0$
\end{itemize}
\end{columns}
\end{frame}

% ============================== SLIDE 9 ==============================
\begin{frame}{Definici\'{o}n alternativa (equivalente)}

\begin{block}{Definici\'{o}n alternativa}
$\{N(t), t \geq 0\}$ es un proceso de Poisson de tasa $\lambda$ si:
\begin{enumerate}
    \item $N(0) = 0$
    \item Tiene incrementos independientes y estacionarios
    \item $N(t) \sim \text{Poisson}(\lambda t)$ para todo $t > 0$
\end{enumerate}
\end{block}

\vspace{0.3cm}
Es decir:
\[
P(N(t) = k) = \frac{e^{-\lambda t}(\lambda t)^k}{k!}, \quad k = 0, 1, 2, \ldots
\]

\vspace{0.3cm}
\textbf{Ejemplo (123):} ?`Probabilidad de exactamente 10 llamadas en 2 minutos?
\[
P(N(2) = 10) = \frac{e^{-8} \cdot 8^{10}}{10!} = \frac{e^{-8} \cdot 1073741824}{3628800} \approx 0.0993
\]
con $\lambda t = 4 \times 2 = 8$.
\end{frame}

% ======================================================================
\section{Distribuci\'{o}n de Poisson}
% ======================================================================

% ============================== SLIDE 10 ==============================
\begin{frame}{Distribuci\'{o}n de Poisson: resumen}

\begin{block}{$X \sim \text{Poisson}(\mu)$}
Modela el \textbf{n\'{u}mero de eventos} en un intervalo fijo.
\[
P(X = k) = \frac{e^{-\mu}\,\mu^k}{k!}, \quad k = 0, 1, 2, \ldots
\]
\end{block}

\vspace{0.2cm}
\textbf{Propiedades:}
\begin{columns}
\column{0.45\textwidth}
\begin{itemize}
    \item $E[X] = \mu$
    \item $\text{Var}(X) = \mu$
    \item $E[X] = \text{Var}(X)$ siempre
    \item Soporte: $\{0, 1, 2, \ldots\}$
\end{itemize}

\column{0.55\textwidth}
\textbf{Propiedad de suma:}

Si $X_1 \sim \text{Poisson}(\mu_1)$ y $X_2 \sim \text{Poisson}(\mu_2)$ son independientes:
\[
X_1 + X_2 \sim \text{Poisson}(\mu_1 + \mu_2)
\]
\end{columns}

\vspace{0.3cm}
\textbf{Ejemplo (123):} En 5 minutos, el n\'{u}mero de llamadas $N(5) \sim \text{Poisson}(20)$.
\[
E[N(5)] = 20, \quad \text{Var}(N(5)) = 20, \quad \text{DE}(N(5)) = \sqrt{20} \approx 4.47
\]
\end{frame}

% ============================== SLIDE 11 ==============================
\begin{frame}{Forma de la distribuci\'{o}n de Poisson}

\begin{center}
\begin{tikzpicture}
\begin{axis}[
    width=13cm, height=6.5cm,
    xlabel={$k$ (n\'{u}mero de eventos)},
    ylabel={$P(X = k)$},
    xmin=-0.5, xmax=20,
    ymin=0, ymax=0.4,
    grid=both, grid style={gray!15},
    legend pos=north east,
    ybar,
    bar width=2pt,
    xtick={0,2,...,20},
]

% mu=2
\addplot[fill=azuloscuro, draw=azuloscuro, opacity=0.7] coordinates {
(0,0.1353) (1,0.2707) (2,0.2707) (3,0.1804) (4,0.0902)
(5,0.0361) (6,0.0120) (7,0.0034) (8,0.0009) (9,0.0002)
};
\addlegendentry{$\mu = 2$}

% mu=5
\addplot[fill=rojoclaro, draw=rojoclaro, opacity=0.7] coordinates {
(0,0.0067) (1,0.0337) (2,0.0842) (3,0.1404) (4,0.1755)
(5,0.1755) (6,0.1462) (7,0.1044) (8,0.0653) (9,0.0363)
(10,0.0181) (11,0.0082) (12,0.0034) (13,0.0013) (14,0.0005)
};
\addlegendentry{$\mu = 5$}

% mu=10
\addplot[fill=verdeoscuro, draw=verdeoscuro, opacity=0.7] coordinates {
(0,0.0000) (1,0.0005) (2,0.0023) (3,0.0076) (4,0.0189)
(5,0.0378) (6,0.0631) (7,0.0901) (8,0.1126) (9,0.1251)
(10,0.1251) (11,0.1137) (12,0.0948) (13,0.0729) (14,0.0521)
(15,0.0347) (16,0.0217) (17,0.0128) (18,0.0071) (19,0.0037) (20,0.0019)
};
\addlegendentry{$\mu = 10$}

\end{axis}
\end{tikzpicture}
\end{center}

A mayor $\mu$, la distribuci\'{o}n se desplaza a la derecha y se vuelve m\'{a}s sim\'{e}trica (aproximaci\'{o}n normal).
\end{frame}

% ============================== SLIDE 12 ==============================
\begin{frame}{Ejemplo completo: Poisson en acci\'{o}n}

\textbf{Situaci\'{o}n:} Un sitio web recibe en promedio $\lambda = 3$ visitas por segundo.

\vspace{0.3cm}
\textbf{Pregunta 1:} ?`Probabilidad de 0 visitas en 1 segundo?
\[
P(N(1) = 0) = \frac{e^{-3}\cdot 3^0}{0!} = e^{-3} \approx 0.0498 \approx 5\%
\]

\textbf{Pregunta 2:} ?`Probabilidad de m\'{a}s de 5 visitas en 1 segundo?
\[
P(N(1) > 5) = 1 - \sum_{k=0}^{5} \frac{e^{-3}\cdot 3^k}{k!} = 1 - 0.9161 = 0.0839 \approx 8.4\%
\]

\textbf{Pregunta 3:} ?`Probabilidad de exactamente 20 visitas en 5 segundos?

Con $\mu = \lambda t = 3\times 5 = 15$:
\[
P(N(5) = 20) = \frac{e^{-15}\cdot 15^{20}}{20!} \approx 0.0418 \approx 4.2\%
\]
\end{frame}

% ======================================================================
\section{Distribuci\'{o}n Exponencial}
% ======================================================================

% ============================== SLIDE 13 ==============================
\begin{frame}{Del conteo al tiempo: distribuci\'{o}n exponencial}

Si los eventos siguen un proceso de Poisson de tasa $\lambda$, ?`cu\'{a}nto tiempo transcurre \textbf{entre eventos consecutivos}?

\vspace{0.3cm}
Sea $T_1$ el tiempo hasta el primer evento. Entonces:
\[
P(T_1 > t) = P(N(t) = 0) = e^{-\lambda t}
\]

Por lo tanto:
\[
F_{T_1}(t) = P(T_1 \leq t) = 1 - e^{-\lambda t}, \quad t \geq 0
\]

\begin{block}{Resultado fundamental}
Los \textbf{tiempos entre llegadas} de un proceso de Poisson($\lambda$) son variables aleatorias \textbf{exponenciales} con par\'{a}metro $\lambda$:
\[
T_i \overset{\text{iid}}{\sim} \text{Exp}(\lambda)
\]
independientes e id\'{e}nticamente distribuidas.
\end{block}

\textbf{Ejemplo (123):} El tiempo entre llamadas es Exp(4), con media $1/4 = 0.25$ minutos = 15 segundos.
\end{frame}

% ============================== SLIDE 14 ==============================
\begin{frame}{Distribuci\'{o}n exponencial: resumen}

\begin{block}{$T \sim \text{Exp}(\lambda)$}
Modela el \textbf{tiempo entre eventos} consecutivos.
\[
f(t) = \lambda e^{-\lambda t}, \quad t \geq 0 \qquad\qquad F(t) = 1 - e^{-\lambda t}
\]
\end{block}

\vspace{0.2cm}
\textbf{Propiedades clave:}
\renewcommand{\arraystretch}{1.4}
\begin{center}
\begin{tabular}{ll}
\toprule
\textbf{Propiedad} & \textbf{Valor} \\
\midrule
Media & $E[T] = 1/\lambda$ \\
Varianza & $\text{Var}(T) = 1/\lambda^2$ \\
Coeficiente de variaci\'{o}n & $CV = 1$ (siempre) \\
Mediana & $\ln(2)/\lambda \approx 0.693/\lambda$ \\
Moda & $0$ \\
\bottomrule
\end{tabular}
\end{center}

\vspace{0.2cm}
\textbf{Ejemplo (123):} $T \sim \text{Exp}(4)$ $\Rightarrow$ $E[T] = 0.25$ min, $\text{Var}(T) = 0.0625$ min$^2$, mediana $= 0.173$ min $\approx 10.4$ seg.
\end{frame}

% ============================== SLIDE 15 ==============================
\begin{frame}{Densidad y CDF de la exponencial}

\begin{columns}
\column{0.5\textwidth}
\begin{center}
\begin{tikzpicture}
\begin{axis}[
    width=7cm, height=5.5cm,
    xlabel={$t$},
    ylabel={$f(t) = \lambda e^{-\lambda t}$},
    title={\textbf{Densidad (PDF)}},
    xmin=0, xmax=3,
    ymin=0, ymax=4.5,
    grid=both, grid style={gray!15},
    legend pos=north east,
    legend style={font=\tiny},
]
\addplot[thick, azuloscuro, domain=0:3, samples=100] {1*exp(-1*x)};
\addlegendentry{$\lambda=1$}
\addplot[thick, rojoclaro, domain=0:3, samples=100] {2*exp(-2*x)};
\addlegendentry{$\lambda=2$}
\addplot[thick, verdeoscuro, domain=0:3, samples=100] {4*exp(-4*x)};
\addlegendentry{$\lambda=4$}
\end{axis}
\end{tikzpicture}
\end{center}

\column{0.5\textwidth}
\begin{center}
\begin{tikzpicture}
\begin{axis}[
    width=7cm, height=5.5cm,
    xlabel={$t$},
    ylabel={$F(t) = 1-e^{-\lambda t}$},
    title={\textbf{Acumulada (CDF)}},
    xmin=0, xmax=3,
    ymin=0, ymax=1.1,
    grid=both, grid style={gray!15},
    legend pos=south east,
    legend style={font=\tiny},
]
\addplot[thick, azuloscuro, domain=0:3, samples=100] {1-exp(-1*x)};
\addlegendentry{$\lambda=1$}
\addplot[thick, rojoclaro, domain=0:3, samples=100] {1-exp(-2*x)};
\addlegendentry{$\lambda=2$}
\addplot[thick, verdeoscuro, domain=0:3, samples=100] {1-exp(-4*x)};
\addlegendentry{$\lambda=4$}
\end{axis}
\end{tikzpicture}
\end{center}
\end{columns}

\vspace{0.3cm}
A mayor $\lambda$, los tiempos entre eventos son m\'{a}s cortos (la densidad se concentra cerca de 0).
\end{frame}

% ============================== SLIDE 16 ==============================
\begin{frame}{Propiedad de falta de memoria}

\begin{block}{Teorema: Falta de memoria (memoryless)}
Si $T \sim \text{Exp}(\lambda)$, entonces para todo $s, t > 0$:
\[
P(T > s + t \mid T > s) = P(T > t)
\]
La exponencial es la \textbf{\'{u}nica} distribuci\'{o}n continua con esta propiedad.
\end{block}

\vspace{0.2cm}
\textbf{Demostraci\'{o}n:}
\[
P(T > s+t \mid T > s) = \frac{P(T > s+t)}{P(T > s)} = \frac{e^{-\lambda(s+t)}}{e^{-\lambda s}} = e^{-\lambda t} = P(T > t) \quad \square
\]

\vspace{0.2cm}
\textbf{Ejemplo (123):} Si han pasado 20 segundos sin llamada, la probabilidad de esperar al menos 10 segundos m\'{a}s es \textbf{la misma} que si acabara de llegar una llamada:
\[
P(T > 30\text{s} \mid T > 20\text{s}) = P(T > 10\text{s}) = e^{-4/6} \approx 0.513
\]
El sistema ``no recuerda'' cu\'{a}nto lleva esperando.
\end{frame}

% ============================== SLIDE 17 ==============================
\begin{frame}{Importancia de la falta de memoria en simulaci\'{o}n}

?`Por qu\'{e} es tan importante esta propiedad?

\vspace{0.3cm}
\begin{columns}
\column{0.5\textwidth}
\textbf{Simplifica el an\'{a}lisis:}
\begin{itemize}
    \item El estado futuro del sistema solo depende del estado \textbf{actual}, no de la historia
    \item Esto hace que el sistema sea \textbf{Markoviano}
    \item Permite soluciones anal\'{i}ticas cerradas ($M/M/c$, etc.)
\end{itemize}

\column{0.5\textwidth}
\textbf{Pero es tambi\'{e}n una limitaci\'{o}n:}
\begin{itemize}
    \item En la realidad, muchos tiempos de servicio \textbf{no} son memoryless
    \item Un m\'{e}dico que lleva 30 min con un paciente probablemente terminar\'{a} pronto
    \item Esto motiva usar distribuciones m\'{a}s generales (Gamma, Weibull, etc.)
\end{itemize}
\end{columns}

\vspace{0.4cm}
\begin{alertblock}{Conexi\'{o}n fundamental}
Proceso de Poisson $\Leftrightarrow$ Tiempos inter-evento exponenciales $\Leftrightarrow$ Propiedad de falta de memoria. Los tres son equivalentes.
\end{alertblock}
\end{frame}

% ============================== SLIDE 18 ==============================
\begin{frame}{Ejemplo: c\'{a}lculos con la exponencial}

\textbf{Situaci\'{o}n:} Un cajero de banco atiende clientes con tiempo de servicio $S \sim \text{Exp}(\mu = 6)$ clientes/hora (media = 10 min).

\vspace{0.3cm}
\textbf{Pregunta 1:} ?`Probabilidad de que un servicio dure m\'{a}s de 15 minutos?
\[
P(S > 0.25\text{ h}) = e^{-6(0.25)} = e^{-1.5} \approx 0.2231 \approx 22.3\%
\]

\textbf{Pregunta 2:} ?`Probabilidad de terminar en menos de 5 minutos?
\[
P(S < 1/12\text{ h}) = 1 - e^{-6/12} = 1 - e^{-0.5} \approx 0.3935 \approx 39.4\%
\]

\textbf{Pregunta 3:} Si ya lleva 10 min, ?`prob.\ de durar al menos 5 min m\'{a}s?
\[
P(S > 15 \mid S > 10) = P(S > 5) = e^{-6(1/12)} = e^{-0.5} \approx 0.6065
\]
(falta de memoria: es como si empezara de nuevo).
\end{frame}

% ======================================================================
\section{Distribuci\'{o}n Gamma (Erlang)}
% ======================================================================

% ============================== SLIDE 19 ==============================
\begin{frame}{Del primer evento al $n$-\'{e}simo: distribuci\'{o}n Gamma}

?`Cu\'{a}nto tiempo $S_n$ pasa hasta que ocurran los primeros $n$ eventos?

\[
S_n = T_1 + T_2 + \cdots + T_n, \quad T_i \overset{\text{iid}}{\sim} \text{Exp}(\lambda)
\]

\begin{block}{Resultado}
$S_n$ sigue una distribuci\'{o}n \textbf{Gamma} (o Erlang si $n$ entero):
\[
S_n \sim \text{Gamma}(n, \lambda)
\]
con densidad:
\[
f_{S_n}(t) = \frac{\lambda^n\, t^{n-1}\, e^{-\lambda t}}{(n-1)!}, \quad t > 0
\]
\end{block}

\vspace{0.2cm}
\textbf{Propiedades:}
\begin{itemize}
    \item $E[S_n] = n/\lambda$
    \item $\text{Var}(S_n) = n/\lambda^2$
    \item $CV(S_n) = 1/\sqrt{n}$ (decrece con $n$)
\end{itemize}
\end{frame}


\section{Distribución Gamma (Erlang)}
% ======================================================================
% ============================== SLIDE 19 ==============================
\begin{frame}{Del primer evento al $n$-ésimo: distribución Gamma}
¿Cuánto tiempo $S_n$ pasa hasta que ocurran los primeros $n$ eventos?
\[
S_n = T_1 + T_2 + \cdots + T_n, \quad T_i \overset{\text{iid}}{\sim} \text{Exp}(\lambda)
\]
\begin{block}{Resultado}
$S_n \sim \text{Gamma}(n, \lambda)$ con densidad y \textbf{acumulada}:
\[
f_{S_n}(t) = \frac{\lambda^n\, t^{n-1}\, e^{-\lambda t}}{(n-1)!}, \quad t > 0
\qquad\qquad
F_{S_n}(t) = P(S_n \le t) = 1 - \sum_{k=0}^{n-1} \frac{e^{-\lambda t}(\lambda t)^k}{k!}
\]
\end{block}

\begin{alertblock}{Dualidad Poisson--Gamma (¡clave!)}
La acumulada de la Gamma \textbf{es} la cola de la Poisson:
\[
\underbrace{P(S_n \le t)}_{\substack{\text{el }n\text{-ésimo evento}\\\text{ocurre antes de }t}}
\;=\;
\underbrace{P\!\left(N(t) \ge n\right)}_{\substack{\text{al menos }n\text{ eventos}\\\text{ocurren en }[0,t]}}
\;=\;
1 - \sum_{k=0}^{n-1} \frac{e^{-\lambda t}(\lambda t)^k}{k!}
\]
\textbf{No se necesitan tablas de la Gamma}: basta sumar términos de Poisson.
\end{alertblock}

\vspace{0.15cm}
\textbf{Propiedades:} \;
$E[S_n] = n/\lambda$, \quad
$\text{Var}(S_n) = n/\lambda^2$, \quad
$CV(S_n) = 1/\sqrt{n}$ (decrece con $n$).
\end{frame}

% ============================== SLIDE 20 ==============================
\begin{frame}{Ejemplo: tiempo hasta el $n$-\'{e}simo evento}

\textbf{Situaci\'{o}n (123):} $\lambda = 4$ llamadas/min. ?`Cu\'{a}nto tiempo hasta la 5\textsuperscript{a} llamada?

\vspace{0.2cm}
$S_5 \sim \text{Gamma}(5, 4)$ con:
\begin{itemize}
    \item $E[S_5] = 5/4 = 1.25$ minutos
    \item $\text{DE}(S_5) = \sqrt{5}/4 \approx 0.559$ minutos
\end{itemize}

\vspace{0.2cm}
\begin{center}
\begin{tikzpicture}
\begin{axis}[
    width=11cm, height=5cm,
    xlabel={$t$ (minutos)},
    ylabel={$f_{S_n}(t)$},
    xmin=0, xmax=5,
    ymin=0, ymax=1.8,
    grid=both, grid style={gray!15},
    legend pos=north east,
    legend style={font=\small},
]
% Gamma(1,4) = Exp(4)
\addplot[thick, azuloscuro, domain=0.01:5, samples=150] {4*exp(-4*x)};
\addlegendentry{$n=1$ (Exp)}

% Gamma(3,4)
\addplot[thick, rojoclaro, domain=0.01:5, samples=150] {(64*x^2*exp(-4*x))/2};
\addlegendentry{$n=3$}

% Gamma(5,4)
\addplot[thick, verdeoscuro, domain=0.01:5, samples=150] {(1024*x^4*exp(-4*x))/24};
\addlegendentry{$n=5$}

% Gamma(10,4)
\addplot[thick, morado, domain=0.01:5, samples=150] {(4^10*x^9*exp(-4*x))/362880};
\addlegendentry{$n=10$}

\end{axis}
\end{tikzpicture}
\end{center}

A mayor $n$, la densidad se concentra m\'{a}s y parece una campana (TLC).
\end{frame}

% ============================== SLIDE 21 ==============================
\begin{frame}{Distribuci\'{o}n Gamma general}

\begin{block}{$X \sim \text{Gamma}(\alpha, \beta)$}
\[
f(x) = \frac{\beta^\alpha}{\Gamma(\alpha)}\, x^{\alpha-1}\, e^{-\beta x}, \quad x > 0
\]
donde $\alpha > 0$ (forma) y $\beta > 0$ (tasa). $\Gamma(\alpha) = \int_0^\infty t^{\alpha-1}e^{-t}\,dt$.
\end{block}

\vspace{0.2cm}
\renewcommand{\arraystretch}{1.4}
\begin{center}
\begin{tabular}{ll}
\toprule
\textbf{Propiedad} & \textbf{Valor} \\
\midrule
Media & $E[X] = \alpha / \beta$ \\
Varianza & $\text{Var}(X) = \alpha / \beta^2$ \\
$CV$ & $1/\sqrt{\alpha}$ \\
Moda & $(\alpha - 1)/\beta$ si $\alpha \geq 1$ \\
\bottomrule
\end{tabular}
\end{center}

\vspace{0.2cm}
\textbf{Casos especiales:}
\begin{itemize}
    \item $\text{Gamma}(1, \lambda) = \text{Exp}(\lambda)$
    \item $\text{Gamma}(n, \lambda)$ con $n \in \mathbb{N}$ = \textbf{Erlang-$n$}
    \item $\text{Gamma}(\nu/2, 1/2) = \chi^2(\nu)$ (Ji-cuadrado)
\end{itemize}
\end{frame}

% ============================== SLIDE 22 ==============================
\begin{frame}{Conexi\'{o}n Poisson -- Gamma: dualidad}

\begin{block}{Teorema de dualidad}
Sea $N(t) \sim \text{Poisson}(\lambda t)$ y $S_n \sim \text{Gamma}(n,\lambda)$. Entonces:
\[
P(N(t) \geq n) = P(S_n \leq t)
\]
\end{block}

\textbf{Interpretaci\'{o}n:} ``Al menos $n$ eventos ocurren antes de $t$'' $\equiv$ ``El $n$-\'{e}simo evento ocurre antes de $t$''.

\vspace{0.3cm}
\textbf{Ejemplo (123):} ?`Prob.\ de al menos 5 llamadas en 2 min?

\vspace{0.1cm}
\textbf{V\'{i}a Poisson:} $P(N(2) \geq 5) = 1 - \sum_{k=0}^{4}\frac{e^{-8}\cdot 8^k}{k!} \approx 1 - 0.0996 = 0.9004$

\vspace{0.1cm}
\textbf{V\'{i}a Gamma:} $P(S_5 \leq 2) = \int_0^2 \frac{4^5 t^4 e^{-4t}}{4!}\,dt = 0.9004$ \quad $\checkmark$

\vspace{0.3cm}
Ambos caminos dan el mismo resultado. Elegimos el que sea m\'{a}s conveniente para el c\'{a}lculo.
\end{frame}

% ======================================================================
\section{Propiedades del Proceso de Poisson}
% ======================================================================

% ============================== SLIDE 23 ==============================
\begin{frame}{Propiedad 1: Superposici\'{o}n (Merging)}

\begin{block}{Teorema de superposici\'{o}n}
Si $N_1(t)$ y $N_2(t)$ son procesos de Poisson independientes con tasas $\lambda_1$ y $\lambda_2$, entonces su uni\'{o}n $N(t) = N_1(t) + N_2(t)$ es un proceso de Poisson con tasa $\lambda_1 + \lambda_2$.
\end{block}

\vspace{0.2cm}
\begin{center}
\begin{tikzpicture}[>=Stealth, thick]
    % Proceso 1
    \draw[->, azuloscuro] (0,1.5) -- (5,1.5);
    \foreach \x in {0.5, 1.3, 2.8, 3.5, 4.3} {
        \draw[azuloscuro, thick] (\x,1.5) -- (\x,2);
        \fill[azuloscuro] (\x,2) circle (2pt);
    }
    \node[left, azuloscuro] at (0,1.75) {$N_1$: $\lambda_1 = 3$};

    % Proceso 2
    \draw[->, rojoclaro] (0,0) -- (5,0);
    \foreach \x in {0.2, 0.9, 1.7, 2.3, 3.1, 4.0, 4.7} {
        \draw[rojoclaro, thick] (\x,0) -- (\x,0.5);
        \fill[rojoclaro] (\x,0.5) circle (2pt);
    }
    \node[left, rojoclaro] at (0,0.25) {$N_2$: $\lambda_2 = 5$};

    % Flecha de merge
    \draw[->, line width=2pt] (5.5,0.75) -- (6.5,0.75);

    % Proceso combinado
    \draw[->, verdeoscuro] (7,-0.3) -- (13.5,-0.3);
    \foreach \x in {7.2, 7.5, 7.9, 8.3, 8.7, 9.1, 9.5, 9.8, 10.3, 10.8, 11.3, 11.7} {
        \draw[verdeoscuro, thick] (\x,-0.3) -- (\x,0.2);
        \fill[verdeoscuro] (\x,0.2) circle (2pt);
    }
    \node[left, verdeoscuro] at (7,-0.05) {$N$: $\lambda = 8$};
\end{tikzpicture}
\end{center}

\vspace{0.2cm}
\textbf{Ejemplo:} Una tienda recibe clientes por dos puertas: puerta A ($\lambda_1 = 3$/h) y puerta B ($\lambda_2 = 5$/h). El flujo total es Poisson con $\lambda = 8$/h.
\end{frame}

% ============================== SLIDE 24 ==============================
\begin{frame}{Ejemplo de superposici\'{o}n: sistema de urgencias}

\textbf{Situaci\'{o}n:} Un hospital recibe pacientes de tres fuentes independientes:
\begin{itemize}
    \item Ambulancias: $\lambda_1 = 2$ pacientes/hora
    \item Caminando: $\lambda_2 = 8$ pacientes/hora
    \item Remisiones: $\lambda_3 = 1.5$ pacientes/hora
\end{itemize}

\vspace{0.3cm}
\textbf{Flujo total:} $\lambda = 2 + 8 + 1.5 = 11.5$ pacientes/hora (Poisson).

\vspace{0.3cm}
\textbf{Pregunta:} ?`Probabilidad de m\'{a}s de 15 pacientes en 1 hora?
\[
P(N(1) > 15) = 1 - \sum_{k=0}^{15} \frac{e^{-11.5}\cdot 11.5^k}{k!} \approx 1 - 0.8914 = 0.1086
\]

\textbf{Pregunta:} ?`Tiempo promedio entre llegadas al sistema total?
\[
E[T] = \frac{1}{\lambda} = \frac{1}{11.5} \approx 5.22 \text{ minutos}
\]
\end{frame}

% ============================== SLIDE 25 ==============================
\begin{frame}{Propiedad 2: Descomposici\'{o}n (Thinning)}

\begin{block}{Teorema de descomposici\'{o}n (Thinning)}
Si $N(t)$ es Poisson($\lambda$) y cada evento se clasifica de forma independiente como tipo~A con prob.\ $p$ o tipo~B con prob.\ $1-p$, entonces:
\begin{itemize}
    \item $N_A(t) \sim \text{Poisson}(\lambda p)$
    \item $N_B(t) \sim \text{Poisson}(\lambda(1-p))$
    \item $N_A(t)$ y $N_B(t)$ son \textbf{independientes}
\end{itemize}
\end{block}

\vspace{0.2cm}
\begin{center}
\begin{tikzpicture}[>=Stealth, thick]
    % Proceso original
    \draw[->] (0,0) -- (5,0);
    \foreach \x/\c in {0.3/A, 0.8/B, 1.2/A, 1.9/B, 2.3/A, 2.7/B, 3.1/A, 3.6/B, 4.0/A, 4.5/B} {
        \draw[black, thick] (\x,0) -- (\x,0.5);
    }
    \node[left] at (0,0.25) {$N$: $\lambda$};

    \draw[->, line width=2pt] (5.5,0.5) -- (6.5,1.2);
    \draw[->, line width=2pt] (5.5,-0.2) -- (6.5,-0.8);
    \node at (6,1) {\small $p$};
    \node at (6,-0.6) {\small $1-p$};

    % Tipo A
    \draw[->, azuloscuro] (7,1) -- (12,1);
    \foreach \x in {7.3, 8.2, 9.3, 10.1, 11.0} {
        \draw[azuloscuro, thick] (\x,1) -- (\x,1.5);
        \fill[azuloscuro] (\x,1.5) circle (2pt);
    }
    \node[right, azuloscuro] at (12,1.25) {$N_A$: $\lambda p$};

    % Tipo B
    \draw[->, rojoclaro] (7,-1) -- (12,-1);
    \foreach \x in {7.8, 8.9, 9.7, 10.6, 11.5} {
        \draw[rojoclaro, thick] (\x,-1) -- (\x,-0.5);
        \fill[rojoclaro] (\x,-0.5) circle (2pt);
    }
    \node[right, rojoclaro] at (12,-0.75) {$N_B$: $\lambda(1-p)$};
\end{tikzpicture}
\end{center}
\end{frame}

% ============================== SLIDE 26 ==============================
\begin{frame}{Ejemplo de descomposici\'{o}n: correo electr\'{o}nico}

\textbf{Situaci\'{o}n:} Un servidor de correo recibe $\lambda = 20$ correos/minuto. Cada correo tiene probabilidad $p = 0.15$ de ser spam (independiente).

\vspace{0.3cm}
\textbf{Por descomposici\'{o}n:}
\begin{itemize}
    \item Spam: Poisson($20 \times 0.15 = 3$ correos/min)
    \item Leg\'{i}timos: Poisson($20 \times 0.85 = 17$ correos/min)
    \item Ambos flujos son \textbf{independientes} entre s\'{i}
\end{itemize}

\vspace{0.3cm}
\textbf{Pregunta:} ?`Prob.\ de 0 spam en 1 minuto?
\[
P(N_{\text{spam}}(1) = 0) = e^{-3} \approx 0.0498 \approx 5\%
\]

\textbf{Pregunta:} ?`Prob.\ de exactamente 2 spam \textbf{y} 15 leg\'{i}timos en 1 minuto?

Por independencia:
\[
P(N_S=2) \times P(N_L=15) = \frac{e^{-3}\cdot 3^2}{2!} \times \frac{e^{-17}\cdot 17^{15}}{15!} \approx 0.2240 \times 0.0847 \approx 0.0190
\]
\end{frame}

% ============================== SLIDE 27 ==============================
\begin{frame}{Propiedad 3: Distribuci\'{o}n condicional de llegadas}

\begin{block}{Teorema (Distribuci\'{o}n uniforme condicional)}
Dado que ocurri\'{o} \textbf{exactamente 1 evento} en $(0,t]$, su instante de llegada $\tau$ se distribuye \textbf{uniformemente} en $(0,t]$:
\[
\tau \mid N(t)=1 \;\sim\; \text{Uniforme}(0,t)
\]
\end{block}

\vspace{0.2cm}
\textbf{Generalizaci\'{o}n:} Dado que $N(t)=n$, los $n$ instantes de llegada $\tau_1 < \tau_2 < \cdots < \tau_n$ se distribuyen como las \textbf{estad\'{i}sticas de orden} de $n$ v.a.\ Uniforme$(0,t)$ independientes.

\vspace{0.3cm}
\textbf{Ejemplo (123):} Si sabemos que hubo exactamente 3 llamadas entre 10:00 y 10:05, entonces sus instantes de llegada son como 3 puntos al azar en ese intervalo de 5 minutos, sin ninguna preferencia por ning\'{u}n subintervalo.

\vspace{0.2cm}
Esta propiedad es fundamental para la \textbf{simulaci\'{o}n}: generar $n$ uniformes y ordenarlas es equivalente a generar $n$ exponenciales acumuladas.
\end{frame}

% ============================== SLIDE 28 ==============================
\begin{frame}{Propiedad 4: Competencia entre exponenciales}

\begin{block}{Teorema (M\'{i}nimo de exponenciales)}
Si $T_1 \sim \text{Exp}(\lambda_1)$ y $T_2 \sim \text{Exp}(\lambda_2)$ son independientes:
\begin{enumerate}
    \item $\min(T_1, T_2) \sim \text{Exp}(\lambda_1 + \lambda_2)$
    \item $P(T_1 < T_2) = \dfrac{\lambda_1}{\lambda_1 + \lambda_2}$
\end{enumerate}
\end{block}

\vspace{0.2cm}
\textbf{Ejemplo: dos servidores.} Un servidor A tiene tiempo de servicio Exp(3) y un servidor B tiene Exp(5). Si ambos empiezan al mismo tiempo:

\vspace{0.2cm}
\begin{itemize}
    \item El primero en terminar lo har\'{a} en tiempo $\sim$ Exp(8), con media $= 1/8 = 7.5$ min.
    \item Probabilidad de que A termine primero: $3/(3+5) = 37.5\%$.
    \item Probabilidad de que B termine primero: $5/(3+5) = 62.5\%$.
\end{itemize}

\vspace{0.2cm}
Se generaliza a $n$ exponenciales: $\min(T_1, \ldots, T_n) \sim \text{Exp}(\sum \lambda_i)$.
\end{frame}

% ============================== SLIDE 29 ==============================
\begin{frame}{Ejemplo: competencia en un sistema de colas}

\textbf{Situaci\'{o}n:} Un sistema $M/M/2$ con $\lambda = 5$, $\mu_1 = 3$, $\mu_2 = 3$. Ambos servidores est\'{a}n ocupados. ?`Qu\'{e} pasa primero: llega un nuevo cliente o termina un servicio?

\vspace{0.3cm}
Tres ``relojes'' exponenciales compitiendo:
\begin{itemize}
    \item Pr\'{o}xima llegada: $T_A \sim \text{Exp}(5)$
    \item Fin servicio 1: $T_{S1} \sim \text{Exp}(3)$
    \item Fin servicio 2: $T_{S2} \sim \text{Exp}(3)$
\end{itemize}

\vspace{0.2cm}
\textbf{El m\'{i}nimo:} $\min(T_A, T_{S1}, T_{S2}) \sim \text{Exp}(5+3+3) = \text{Exp}(11)$.

\vspace{0.2cm}
\textbf{Probabilidades:}
\begin{align}
P(\text{llega cliente primero}) &= \frac{5}{11} \approx 45.5\% \\
P(\text{servidor 1 termina primero}) &= \frac{3}{11} \approx 27.3\% \\
P(\text{servidor 2 termina primero}) &= \frac{3}{11} \approx 27.3\%
\end{align}

Esta es exactamente la l\'{o}gica que usa el simulador de eventos discretos para decidir cu\'{a}l evento ocurre primero.
\end{frame}

% ======================================================================
\section{Distribuci\'{o}n Uniforme}
% ======================================================================

% ============================== SLIDE 30 ==============================
\begin{frame}{Distribuci\'{o}n uniforme: complemento esencial}

\begin{block}{$X \sim \text{Uniforme}(a, b)$}
Todos los valores en $[a,b]$ son igualmente probables.
\[
f(x) = \frac{1}{b-a}, \quad a \leq x \leq b \qquad\qquad F(x) = \frac{x-a}{b-a}
\]
\end{block}

\vspace{0.2cm}
\renewcommand{\arraystretch}{1.4}
\begin{center}
\begin{tabular}{ll}
\toprule
\textbf{Propiedad} & \textbf{Valor} \\
\midrule
Media & $E[X] = (a+b)/2$ \\
Varianza & $\text{Var}(X) = (b-a)^2/12$ \\
Soporte & $[a, b]$ \\
\bottomrule
\end{tabular}
\end{center}

\vspace{0.3cm}
\textbf{Rol en simulaci\'{o}n:}
\begin{itemize}
    \item $U \sim \text{Uniforme}(0,1)$ es la \textbf{base} de toda generaci\'{o}n de n\'{u}meros aleatorios
    \item M\'{e}todo de la transformada inversa: $X = F^{-1}(U)$
    \item Para la exponencial: $T = -\frac{1}{\lambda}\ln(U) \sim \text{Exp}(\lambda)$
\end{itemize}
\end{frame}

% ============================== SLIDE 31 ==============================
\begin{frame}{Ejemplo: generaci\'{o}n de tiempos exponenciales}

\textbf{Transformada inversa para Exp($\lambda = 4$):}

\vspace{0.2cm}
\begin{center}
\begin{tikzpicture}
\begin{axis}[
    width=12cm, height=5.5cm,
    xlabel={$U \sim \text{Uniforme}(0,1)$},
    ylabel={$T = -\frac{1}{4}\ln(U)$},
    xmin=0, xmax=1,
    ymin=0, ymax=2,
    grid=both, grid style={gray!15},
    legend pos=north east,
]
\addplot[thick, azuloscuro, domain=0.001:0.999, samples=200] {-0.25*ln(x)};
\addlegendentry{$F^{-1}(u) = -\frac{1}{\lambda}\ln(u)$}

% Ejemplo
\draw[dashed, rojoclaro, thick] (axis cs:0.3,0) -- (axis cs:0.3,0.301);
\draw[dashed, rojoclaro, thick] (axis cs:0,0.301) -- (axis cs:0.3,0.301);
\node[left, rojoclaro, font=\small] at (axis cs:0,0.301) {$0.301$};
\node[below, rojoclaro, font=\small] at (axis cs:0.3,0) {$0.30$};

\draw[dashed, verdeoscuro, thick] (axis cs:0.7,0) -- (axis cs:0.7,0.0892);
\draw[dashed, verdeoscuro, thick] (axis cs:0,0.0892) -- (axis cs:0.7,0.0892);
\node[left, verdeoscuro, font=\small] at (axis cs:0,0.0892) {$0.089$};
\node[below, verdeoscuro, font=\small] at (axis cs:0.7,0) {$0.70$};

\end{axis}
\end{tikzpicture}
\end{center}

\vspace{0.1cm}
Si $U = 0.30$: $T = -\frac{1}{4}\ln(0.30) = 0.301$ min. Si $U = 0.70$: $T = -\frac{1}{4}\ln(0.70) = 0.089$ min.

As\'{i} se generan las llegadas en un simulador del proceso de Poisson.
\end{frame}

% ======================================================================
\section{Proceso de Poisson No Homog\'{e}neo}
% ======================================================================

% ============================== SLIDE 32 ==============================
\begin{frame}{?`Y si la tasa cambia con el tiempo?}

En la realidad, las tasas de llegada rara vez son constantes:

\begin{itemize}
    \item Las llamadas al 123 aumentan en horas pico
    \item M\'{a}s clientes llegan a un restaurante al mediod\'{i}a
    \item M\'{a}s accidentes ocurren en hora de mayor tr\'{a}fico
\end{itemize}

\vspace{0.3cm}
\begin{block}{Proceso de Poisson no homog\'{e}neo (NHPP)}
Un proceso de conteo $\{N(t)\}$ donde la tasa var\'{i}a en el tiempo: $\lambda(t)$.
\[
P(N(t+h) - N(t) = 1) = \lambda(t)\,h + o(h)
\]
El n\'{u}mero de eventos en $(s,t]$ sigue:
\[
N(t) - N(s) \sim \text{Poisson}\!\left(\int_s^t \lambda(u)\,du\right)
\]
\end{block}

Mantiene incrementos independientes pero \textbf{ya no} estacionarios.
\end{frame}

% ============================== SLIDE 33 ==============================
\begin{frame}{Ejemplo: NHPP en un call center}

\textbf{Situaci\'{o}n:} Tasa de llamadas al 123 a lo largo del d\'{i}a:
\[
\lambda(t) = 2 + 4\sin\!\left(\frac{\pi t}{12}\right), \quad 0 \leq t \leq 24 \text{ horas}
\]

\begin{center}
\begin{tikzpicture}
\begin{axis}[
    width=12cm, height=5cm,
    xlabel={Hora del d\'{i}a},
    ylabel={$\lambda(t)$ (llamadas/min)},
    xmin=0, xmax=24,
    ymin=0, ymax=7,
    grid=both, grid style={gray!15},
    xtick={0,3,...,24},
    xticklabels={0:00,3:00,6:00,9:00,12:00,15:00,18:00,21:00,24:00},
    xticklabel style={font=\tiny, rotate=30, anchor=east},
]
\addplot[thick, azuloscuro, domain=0:24, samples=200] {2 + 4*sin(deg(3.14159*x/12))};

% Pico
\draw[->, thick, rojoclaro] (axis cs:6,6.5) -- (axis cs:6,5.8);
\node[above, rojoclaro, font=\small] at (axis cs:6,6.5) {Pico: 6 llam/min};

% Valle
\draw[->, thick, verdeoscuro] (axis cs:18,1) -- (axis cs:18,1.8);
\node[below, verdeoscuro, font=\small] at (axis cs:18,1) {Valle};

% Area sombreada 9am-12pm
\addplot[fill=blue!15, draw=none, domain=9:15] {2 + 4*sin(deg(3.14159*x/12))} \closedcycle;

\end{axis}
\end{tikzpicture}
\end{center}

\vspace{0.1cm}
Llamadas esperadas entre 9:00 y 15:00: $\int_9^{15}\lambda(t)\,dt \approx 26.3$ llamadas por hora $\times$ duraci\'{o}n.
\end{frame}

% ============================== SLIDE 34 ==============================
\begin{frame}{Simulaci\'{o}n del NHPP: m\'{e}todo de thinning}

\textbf{Algoritmo de Lewis-Shedler (thinning):}

\vspace{0.2cm}
Sea $\lambda^* = \max_{t}\lambda(t)$ la tasa m\'{a}xima.

\vspace{0.2cm}
\begin{enumerate}
    \item Generar candidato del Poisson homog\'{e}neo con tasa $\lambda^*$: $t_{\text{cand}} = t + \text{Exp}(\lambda^*)$
    \item Generar $U \sim \text{Uniforme}(0,1)$
    \item \textbf{Aceptar} el evento si $U \leq \lambda(t_{\text{cand}})/\lambda^*$; si no, \textbf{rechazar}
    \item Repetir
\end{enumerate}

\vspace{0.3cm}
\begin{center}
\begin{tikzpicture}[>=Stealth]
    \draw[->] (0,0) -- (10,0) node[right] {$t$};
    \draw[->] (0,0) -- (0,2.5) node[above] {$\lambda$};

    % lambda*
    \draw[dashed, rojoclaro, thick] (0,2) -- (10,2) node[right] {$\lambda^*$};

    % lambda(t)
    \draw[thick, azuloscuro, domain=0:10, samples=100] plot (\x, {1 + sin(\x*60)*0.8});
    \node[azuloscuro, right] at (10,1) {$\lambda(t)$};

    % Candidatos
    \foreach \x/\a in {1.2/1, 2.5/0, 3.8/1, 5.0/1, 6.1/0, 7.3/1, 8.5/1, 9.2/0} {
        \ifnum\a=1
            \fill[verdeoscuro] (\x,0) circle (3pt);
            \node[below, verdeoscuro, font=\tiny] at (\x,-0.15) {\checkmark};
        \else
            \draw[rojoclaro, thick] (\x,-0.1) -- ++(0.15,0.2) -- ++(-0.3,0) -- cycle;
            \node[below, rojoclaro, font=\tiny] at (\x,-0.15) {$\times$};
        \fi
    }
\end{tikzpicture}
\end{center}
\end{frame}

% ======================================================================
\section{Otras Distribuciones Clave}
% ======================================================================

% ============================== SLIDE 35 ==============================
\begin{frame}{Distribuci\'{o}n Normal (Gaussiana)}

\begin{block}{$X \sim N(\mu, \sigma^2)$}
\[
f(x) = \frac{1}{\sigma\sqrt{2\pi}} \exp\!\left(-\frac{(x-\mu)^2}{2\sigma^2}\right), \quad -\infty < x < \infty
\]
\end{block}

\vspace{0.1cm}
\begin{columns}
\column{0.5\textwidth}
\textbf{Propiedades:}
\begin{itemize}
    \item $E[X] = \mu$, $\text{Var}(X) = \sigma^2$
    \item Sim\'{e}trica alrededor de $\mu$
    \item Regla 68-95-99.7
    \item Cerrada bajo combinaciones lineales
\end{itemize}

\vspace{0.2cm}
\textbf{Rol en simulaci\'{o}n:}
\begin{itemize}
    \item TLC: promedios de muchas v.a.\ son aprox.\ normales
    \item Intervalos de confianza
    \item Tiempos de proceso con poca variabilidad
\end{itemize}

\column{0.5\textwidth}
\begin{center}
\begin{tikzpicture}
\begin{axis}[
    width=6.5cm, height=5cm,
    xlabel={$x$},
    ylabel={$f(x)$},
    xmin=-4, xmax=8,
    ymin=0, ymax=0.45,
    grid=both, grid style={gray!15},
    legend pos=north east,
    legend style={font=\tiny},
]
\addplot[thick, azuloscuro, domain=-4:8, samples=100] {exp(-x^2/2)/sqrt(2*pi)};
\addlegendentry{$N(0,1)$}
\addplot[thick, rojoclaro, domain=-4:8, samples=100] {exp(-(x-2)^2/8)/sqrt(8*pi)};
\addlegendentry{$N(2,4)$}
\addplot[thick, verdeoscuro, domain=-4:8, samples=100] {exp(-(x-3)^2/2)/sqrt(2*pi)};
\addlegendentry{$N(3,1)$}
\end{axis}
\end{tikzpicture}
\end{center}
\end{columns}
\end{frame}

% ============================== SLIDE 36 ==============================
\begin{frame}{Ejemplo: Normal y TLC en simulaci\'{o}n}

\textbf{Situaci\'{o}n:} Estimamos $\bar{W}$ (tiempo promedio en sistema) con $n=100$ observaciones. Cada $W_i$ tiene media $\mu_W = 1.09$ min y DE $= 0.8$ min.

\vspace{0.3cm}
Por el \textbf{Teorema Central del L\'{i}mite}:
\[
\bar{W} \approx N\!\left(\mu_W,\; \frac{\sigma^2}{n}\right) = N\!\left(1.09,\; \frac{0.64}{100}\right)
\]

\vspace{0.2cm}
\textbf{Intervalo de confianza al 95\%:}
\[
\bar{W} \pm z_{0.025}\cdot\frac{\sigma}{\sqrt{n}} = 1.09 \pm 1.96 \cdot \frac{0.8}{\sqrt{100}} = 1.09 \pm 0.157
\]
\[
\text{IC: } [0.933, \; 1.247]
\]

\vspace{0.2cm}
Con $n = 400$: IC $= 1.09 \pm 0.078 = [1.012, 1.168]$ (m\'{a}s estrecho).

\vspace{0.2cm}
La normal es esencial para cuantificar la \textbf{precisi\'{o}n} de nuestras estimaciones por simulaci\'{o}n.
\end{frame}

% ============================== SLIDE 37 ==============================
\begin{frame}{Distribuci\'{o}n Weibull}

\begin{block}{$X \sim \text{Weibull}(\alpha, \beta)$}
\[
f(x) = \frac{\alpha}{\beta}\left(\frac{x}{\beta}\right)^{\alpha-1}\exp\!\left(-\left(\frac{x}{\beta}\right)^{\alpha}\right), \quad x > 0
\]
$\alpha$ = forma, $\beta$ = escala.
\end{block}

\vspace{0.2cm}
\textbf{Propiedades especiales:}
\begin{itemize}
    \item $\alpha = 1$: Weibull $= $ Exponencial (tasa constante de falla)
    \item $\alpha < 1$: Tasa de falla \textbf{decreciente} (mortalidad infantil)
    \item $\alpha > 1$: Tasa de falla \textbf{creciente} (desgaste, envejecimiento)
\end{itemize}

\vspace{0.3cm}
\textbf{Ejemplo:} Vida \'{u}til de un motor con $\alpha = 2.5$, $\beta = 5000$ horas.
\begin{itemize}
    \item $E[X] = 5000\,\Gamma(1 + 1/2.5) \approx 5000 \times 0.887 = 4435$ horas
    \item La tasa de falla \textbf{aumenta} con el tiempo (desgaste), a diferencia de la exponencial
\end{itemize}
\end{frame}

% ============================== SLIDE 38 ==============================
\begin{frame}{Distribuci\'{o}n Lognormal}

\begin{block}{$X \sim \text{Lognormal}(\mu, \sigma^2)$}
Si $\ln(X) \sim N(\mu, \sigma^2)$, entonces $X$ es lognormal.
\[
f(x) = \frac{1}{x\sigma\sqrt{2\pi}}\exp\!\left(-\frac{(\ln x - \mu)^2}{2\sigma^2}\right), \quad x > 0
\]
\end{block}

\vspace{0.2cm}
\textbf{Propiedades:}
\begin{itemize}
    \item $E[X] = e^{\mu + \sigma^2/2}$
    \item $\text{Var}(X) = e^{2\mu+\sigma^2}(e^{\sigma^2}-1)$
    \item Asim\'{e}trica a la derecha (cola pesada)
    \item Siempre positiva
\end{itemize}

\vspace{0.3cm}
\textbf{Ejemplo:} Tiempos de reparaci\'{o}n de m\'{a}quinas, con $\mu = 1.5$, $\sigma = 0.8$:
\begin{itemize}
    \item $E[X] = e^{1.5 + 0.32} = e^{1.82} \approx 6.17$ horas
    \item La mayor\'{i}a de reparaciones son r\'{a}pidas, pero ocasionalmente hay reparaciones muy largas
\end{itemize}
\end{frame}

% ======================================================================
\section{Resumen de Distribuciones}
% ======================================================================

% ============================== SLIDE 39 ==============================
\begin{frame}{Mapa de distribuciones en simulaci\'{o}n de colas}

\begin{center}
\begin{tikzpicture}[>=Stealth, thick,
    cbox/.style={draw, rounded corners, minimum width=2.5cm, minimum height=0.9cm, align=center, font=\small\bfseries},
]
    % Poisson
    \node[cbox, fill=blue!20] (poi) at (0,3) {Poisson($\lambda t$)};
    \node[below, font=\tiny, text width=3cm, align=center] at (0,2.4) {N\'{u}mero de eventos\\en intervalo fijo};

    % Exponencial
    \node[cbox, fill=red!20] (exp) at (-5,0) {Exponencial($\lambda$)};
    \node[below, font=\tiny, text width=3cm, align=center] at (-5,-0.6) {Tiempo entre\\eventos consecutivos};

    % Gamma
    \node[cbox, fill=green!20] (gam) at (0,0) {Gamma($n, \lambda$)};
    \node[below, font=\tiny, text width=3cm, align=center] at (0,-0.6) {Tiempo hasta el\\$n$-\'{e}simo evento};

    % Uniforme
    \node[cbox, fill=orange!20] (uni) at (5,0) {Uniforme$(0,t)$};
    \node[below, font=\tiny, text width=3cm, align=center] at (5,-0.6) {Instantes de llegada\\condicionados};

    % Normal
    \node[cbox, fill=purple!20] (nor) at (-5,-3) {Normal($\mu,\sigma^2$)};
    \node[below, font=\tiny, text width=3cm, align=center] at (-5,-3.6) {Promedios, IC\\(TLC)};

    % Weibull
    \node[cbox, fill=yellow!30] (wei) at (0,-3) {Weibull($\alpha,\beta$)};
    \node[below, font=\tiny, text width=3cm, align=center] at (0,-3.6) {Tiempos de vida\\con desgaste};

    % Lognormal
    \node[cbox, fill=cyan!20] (logn) at (5,-3) {Lognormal($\mu,\sigma^2$)};
    \node[below, font=\tiny, text width=3cm, align=center] at (5,-3.6) {Tiempos de servicio\\con cola pesada};

    % Flechas
    \draw[<->, thick] (poi) -- (exp) node[midway, above left, font=\tiny] {$T \sim$ Exp $\Leftrightarrow$ $N \sim$ Poi};
    \draw[->, thick] (exp) -- (gam) node[midway, above, font=\tiny] {Suma de $n$};
    \draw[->, thick] (poi) -- (uni) node[midway, above right, font=\tiny] {$N(t)=n$ dado};
    \draw[->, thick] (poi) -- (gam) node[midway, right, font=\tiny] {Dualidad};
    \draw[->, thick, dashed] (exp) -- (wei) node[midway, left, font=\tiny] {$\alpha=1$};
    \draw[->, thick, dashed] (gam) -- (nor) node[midway, above, font=\tiny] {$n$ grande (TLC)};

\end{tikzpicture}
\end{center}
\end{frame}

% ============================== SLIDE 40 ==============================
\begin{frame}{Tabla resumen de distribuciones}

\begin{center}
\renewcommand{\arraystretch}{1.3}
\small
\begin{tabular}{lcccc}
\toprule
\textbf{Distribuci\'{o}n} & \textbf{Soporte} & $E[X]$ & $\text{Var}(X)$ & \textbf{Uso t\'{i}pico} \\
\midrule
Poisson($\mu$) & $\{0,1,2,\ldots\}$ & $\mu$ & $\mu$ & Conteo de eventos \\
Exp($\lambda$) & $(0,\infty)$ & $1/\lambda$ & $1/\lambda^2$ & Tiempos entre llegadas \\
Gamma($\alpha,\beta$) & $(0,\infty)$ & $\alpha/\beta$ & $\alpha/\beta^2$ & Suma de exponenciales \\
Uniforme$(a,b)$ & $[a,b]$ & $\frac{a+b}{2}$ & $\frac{(b-a)^2}{12}$ & Generaci\'{o}n aleatoria \\
Normal$(\mu,\sigma^2)$ & $\mathbb{R}$ & $\mu$ & $\sigma^2$ & Promedios (TLC) \\
Weibull$(\alpha,\beta)$ & $(0,\infty)$ & $\beta\Gamma(1\!+\!1/\alpha)$ & --- & Confiabilidad \\
Lognormal$(\mu,\sigma^2)$ & $(0,\infty)$ & $e^{\mu+\sigma^2/2}$ & --- & Tiempos de reparaci\'{o}n \\
\bottomrule
\end{tabular}
\end{center}

\vspace{0.3cm}
\begin{alertblock}{Nota para la simulaci\'{o}n}
La notaci\'{o}n de Kendall $A/B/c$ usa: $M$ = exponencial (Markoviano), $D$ = determin\'{i}stico, $G$ = general, $E_k$ = Erlang-$k$. El proceso de Poisson da la ``$M$'' en las llegadas y servicios.
\end{alertblock}
\end{frame}

% ======================================================================
\section{Simulaci\'{o}n del Proceso de Poisson}
% ======================================================================

% ============================== SLIDE 41 ==============================
\begin{frame}{C\'{o}mo simular un proceso de Poisson homog\'{e}neo}

\textbf{M\'{e}todo 1: V\'{i}a tiempos entre llegadas}

\vspace{0.2cm}
\begin{enumerate}
    \item Inicializar: $t \leftarrow 0$, $n \leftarrow 0$
    \item Generar $U \sim \text{Uniforme}(0,1)$
    \item Calcular $T = -\frac{1}{\lambda}\ln(U)$
    \item Actualizar: $t \leftarrow t + T$, $n \leftarrow n + 1$
    \item Registrar evento en el tiempo $t$
    \item Repetir desde paso 2 hasta alcanzar el horizonte deseado
\end{enumerate}

\vspace{0.3cm}
\textbf{Ejemplo con $\lambda = 4$:}

\renewcommand{\arraystretch}{1.2}
\begin{center}
\small
\begin{tabular}{ccccc}
\toprule
$n$ & $U$ & $T = -\frac{1}{4}\ln(U)$ & Llegada $t$ & $N(t)$ \\
\midrule
1 & 0.372 & 0.247 & 0.247 & 1 \\
2 & 0.891 & 0.029 & 0.276 & 2 \\
3 & 0.156 & 0.464 & 0.740 & 3 \\
4 & 0.625 & 0.118 & 0.858 & 4 \\
5 & 0.044 & 0.781 & 1.639 & 5 \\
\bottomrule
\end{tabular}
\end{center}
\end{frame}

% ============================== SLIDE 42 ==============================
\begin{frame}{M\'{e}todo 2: V\'{i}a conteo directo}

\textbf{Para generar $N(t)$ directamente en un intervalo $[0, T]$:}

\vspace{0.2cm}
\begin{enumerate}
    \item Generar $N \sim \text{Poisson}(\lambda T)$
    \item Generar $N$ uniformes independientes $U_1, \ldots, U_N \sim \text{Uniforme}(0, T)$
    \item Ordenar: $U_{(1)} < U_{(2)} < \cdots < U_{(N)}$
    \item Los $U_{(i)}$ son los instantes de llegada
\end{enumerate}

\vspace{0.3cm}
\textbf{Ejemplo con $\lambda = 4$, $T = 2$ min:}

\vspace{0.2cm}
Paso 1: $N \sim \text{Poisson}(8)$, sale $N = 7$

Paso 2: Generamos 7 uniformes en $[0, 2]$:
\[
0.34, \; 1.72, \; 0.89, \; 0.12, \; 1.45, \; 0.56, \; 1.91
\]

Paso 3: Ordenamos:
\[
0.12, \; 0.34, \; 0.56, \; 0.89, \; 1.45, \; 1.72, \; 1.91
\]

Estos son los 7 instantes de llegada en $[0, 2]$.

Este m\'{e}todo es \'{u}til cuando se necesita el n\'{u}mero total primero (e.g., para asignaci\'{o}n de recursos).
\end{frame}

% ======================================================================
\section{Ejercicios y Aplicaciones}
% ======================================================================

% ============================== SLIDE 43 ==============================
\begin{frame}{Ejercicio integrador 1: aeropuerto}

\textbf{Problema:} Un aeropuerto peque\~{n}o tiene una sola pista. Los aviones llegan seg\'{u}n un proceso de Poisson con $\lambda = 8$ aviones/hora. Cada aterrizaje toma un tiempo Exp($\mu = 12$) por hora.

\vspace{0.3cm}
\textbf{Resuelva:}
\begin{enumerate}
    \item ?`Cu\'{a}l es la probabilidad de que lleguen m\'{a}s de 3 aviones en 15 minutos?

    \textit{Sol:} $N(0.25) \sim \text{Poisson}(2)$. $P(N>3) = 1 - \sum_{k=0}^{3} \frac{e^{-2}\cdot 2^k}{k!} = 1 - 0.8571 = 0.143$.

    \vspace{0.15cm}
    \item ?`Tiempo medio entre llegadas?

    \textit{Sol:} $E[T] = 1/8 = 0.125$ h $= 7.5$ minutos.

    \vspace{0.15cm}
    \item ?`Prob.\ de que un aterrizaje dure m\'{a}s de 10 min?

    \textit{Sol:} $P(S > 1/6) = e^{-12/6} = e^{-2} \approx 0.135$.

    \vspace{0.15cm}
    \item ?`Intensidad de tr\'{a}fico?

    \textit{Sol:} $\rho = \lambda/\mu = 8/12 = 2/3 \approx 0.667$. Sistema estable.
\end{enumerate}
\end{frame}

% ============================== SLIDE 44 ==============================
\begin{frame}{Ejercicio integrador 2: supermercado}

\textbf{Problema:} Un supermercado tiene dos cajas. Los clientes llegan Poisson($\lambda = 10$/h). El 40\% va a la caja r\'{a}pida y el 60\% a la regular.

\vspace{0.3cm}
\textbf{Resuelva usando descomposici\'{o}n:}
\begin{enumerate}
    \item Llegadas a la caja r\'{a}pida: Poisson($10 \times 0.4 = 4$/h)
    \item Llegadas a la caja regular: Poisson($10 \times 0.6 = 6$/h)
    \item Son procesos \textbf{independientes}
\end{enumerate}

\vspace{0.3cm}
\textbf{Preguntas:}
\begin{itemize}
    \item ?`Prob.\ de 0 clientes en caja r\'{a}pida en 30 min? $P(N_R(0.5)=0) = e^{-2} \approx 0.135$

    \item ?`Tiempo medio entre clientes en caja regular? $1/6$ h $= 10$ min

    \item ?`Prob.\ de que lleguen exactamente 5 clientes en total a ambas cajas en 30 min?

    $N(0.5) \sim \text{Poisson}(5)$. $P(N=5) = \frac{e^{-5}\cdot 5^5}{5!} = \frac{e^{-5}\cdot 3125}{120} \approx 0.1755$
\end{itemize}
\end{frame}

% ============================== SLIDE 45 ==============================
\begin{frame}{Ejercicio integrador 3: competencia de exponenciales}

\textbf{Problema:} Tres m\'{a}quinas funcionan en paralelo. Sus tiempos de vida son exponenciales independientes con tasas $\lambda_1 = 0.01$, $\lambda_2 = 0.02$, $\lambda_3 = 0.015$ fallas/hora.

\vspace{0.3cm}
\textbf{a) ?`Cu\'{a}nto tiempo hasta la primera falla?}
\[
\min(T_1,T_2,T_3) \sim \text{Exp}(0.01+0.02+0.015) = \text{Exp}(0.045)
\]
\[
E[\text{primera falla}] = 1/0.045 = 22.2 \text{ horas}
\]

\textbf{b) ?`Qu\'{e} m\'{a}quina tiene mayor probabilidad de fallar primero?}
\[
P(T_2 = \min) = \frac{0.02}{0.045} = 44.4\% \quad \text{(la m\'{a}quina 2, la menos confiable)}
\]

\textbf{c) ?`Prob.\ de que la m\'{a}quina 1 dure m\'{a}s de 50 horas?}
\[
P(T_1 > 50) = e^{-0.01 \times 50} = e^{-0.5} \approx 0.607
\]
\end{frame}

% ======================================================================
\section{Conclusiones}
% ======================================================================

% ============================== SLIDE 46 ==============================
\begin{frame}{Conclusiones principales}

\begin{enumerate}
    \item El \textbf{proceso de Poisson} modela eventos aleatorios con tasa constante $\lambda$, caracterizado por incrementos independientes y estacionarios.

    \vspace{0.15cm}
    \item La \textbf{distribuci\'{o}n de Poisson} cuenta eventos en intervalos fijos; la \textbf{exponencial} mide tiempos entre eventos; la \textbf{Gamma} mide el tiempo hasta el $n$-\'{e}simo evento.

    \vspace{0.15cm}
    \item La \textbf{falta de memoria} de la exponencial es lo que hace ``Markovianos'' a los modelos $M/M/c$.

    \vspace{0.15cm}
    \item Las propiedades de \textbf{superposici\'{o}n} y \textbf{descomposici\'{o}n} permiten analizar sistemas complejos con m\'{u}ltiples fuentes.

    \vspace{0.15cm}
    \item La \textbf{competencia de exponenciales} es el mecanismo central del simulador de eventos discretos.

    \vspace{0.15cm}
    \item Distribuciones como \textbf{Weibull}, \textbf{Lognormal} y \textbf{Normal} complementan el toolkit para modelar fen\'{o}menos m\'{a}s realistas.
\end{enumerate}
\end{frame}

% ============================== SLIDE 47 ==============================
\begin{frame}{Referencias}

\begin{thebibliography}{9}
\small

\bibitem{ross} Ross, S.M. (2014). \textit{Introduction to Probability Models}, 11th ed. Academic Press. Cap.~5: El proceso de Poisson.

\bibitem{law} Law, A.M. (2015). \textit{Simulation Modeling and Analysis}, 5th ed. McGraw-Hill. Cap.~8: Generaci\'{o}n de variables aleatorias.

\bibitem{banks} Banks, J. et al. (2014). \textit{Discrete-Event System Simulation}, 5th ed. Pearson. Cap.~5--6.

\bibitem{gross} Gross, D. et al. (2008). \textit{Fundamentals of Queueing Theory}, 4th ed. Wiley.

\bibitem{cinlar} \c{C}inlar, E. (2013). \textit{Introduction to Stochastic Processes}. Dover. Cap.~4: Procesos de Poisson.

\bibitem{kingman} Kingman, J.F.C. (1993). \textit{Poisson Processes}. Oxford University Press.

\end{thebibliography}
\end{frame}

% ============================== SLIDE 48 ==============================
\begin{frame}{}

\begin{center}
\vspace{1cm}
{\Huge\textbf{?`Preguntas?}}


\vspace{0.5cm}
{\large\textit{``The Poisson process is the most important stochastic process\\in applied probability.''} --- Sheldon M. Ross}


\end{center}
\end{frame}

\end{document}
